\documentclass[10pt]{article}
\usepackage[a4paper,margin=25mm]{geometry}
\usepackage{amsmath,amssymb,amsthm,booktabs}
\newtheorem{theorem}{Theorem}
\title{Coordinates and residues in the Higman--Sims catalogue}
\author{Private L091 component}\date{}
\begin{document}\maketitle
Let $G=\mathrm{HS}$. We use the supplement dated 6 October 2003,
with 23 pages, as the fixed source for the diagram numbers below.
Its statements and notation remain attributed to the original authors.
The complete classification, the local numerical data, the bare
isomorphism distinctions, and the truncation certificates are acquired
from their separately reviewed components. This application identifies
the source coordinates and verifies the small named residues. It does
not replace those earlier proofs.

\begin{theorem}
The numerical profiles of the 86 diagrams of ranks two through six,
with the six explicit amendments below, give a bijection to the 86
classes in the acquired classification. The thirteen printed truncation
assertions have actual incidence isomorphisms. The thirteen occurrences
of a Petersen residue and the two occurrences of an affine-plane residue
have explicit isomorphisms to the Petersen incidence structure and
$\mathrm{AG}(2,4)$, respectively.
\end{theorem}

\paragraph{Coordinates and the two ambiguities.}
For each source diagram, first compare the multiset of type sizes with
every acquired class. Apart from two pairs, this identifies a unique
class. The entry $44320$ in diagram 2.6 is retained in the transcription
but replaced by $44352$ for this comparison: the printed integer does
not divide $|G|=44352000$, so it cannot be an index for the asserted
transitive action.

The two rank-three diagrams 3.22 and 3.23 have identical type sizes and
local numerical data. Their candidate classes are 14 and 15. The acquired
full automorphism proof gives $\operatorname{Aut}(\Gamma_{14})=G$ and
$\operatorname{Aut}(\Gamma_{15})=\operatorname{Aut}(G)$, of respective
orders $44352000$ and $88704000$. These are full typed groups, rather
than merely normalizers of a specified action. Thus the source's two
different automorphism assertions select 14 for 3.22 and 15 for 3.23.
For the remaining tied pair, 5.2 and 5.3, the complete diagram has
respectively one and two quadrangle bonds. Every permutation among
equal-sized types is checked. This selects classes 83 and 84.

The finite coordinate list records a permutation from each source's
ordered nodes to the actual model's types. The resulting 86 class
indices are precisely $1,\ldots,86$, each once. This identifies classes
by proved invariants within a complete classification; it makes no
claim that an unavailable historical generator list is literally the
same list of permutations as ours.

\paragraph{Literal discrepancies.}
An ordered triple $(a,g,b)$ records the eccentricities of the two
specified vertex types and half the girth of their incidence graph.
The finite record retains both the printed and the actual values.
The six amendments are as follows.
\begin{center}\begin{tabular}{lp{10.6cm}}
\toprule
Diagram & Amendment\\\midrule
2.5 & The node order at type $36960$ is $15$, in place of $11$.\\
2.6 & The second type has size $44352$, in place of $44320$.\\
3.15 & The right arm, ordered from the type of size $352$ to that of size
$1100$, has triple $(5,4,6)$, in place of $(6,4,5)$.\\
3.30 & The top node order is $3$, in place of $1$.\\
4.22 & The bottom bond is a digon; the omitted diagonal joining the
two types of size $352$ has a single bond.\\
4.28 & The bond from type $352$ to type $1386000$ has triple $(7,5,8)$;
the bond from type $1386000$ to type $1100$ has triple $(4,3,3)$ and
is a complete-graph incidence structure.\\
\bottomrule
\end{tabular}\end{center}

\paragraph{Why these are amendments.}
For each of these six rows the type-size multiset, after the explicitly
declared index amendment in 2.6, selects a unique class. All compatible
type permutations are tested against the acquired local data. Therefore
the stated mismatch cannot be removed by choosing another class or by
silently exchanging equal-sized types. In particular, both possible
exchanges of the $352$ types in 4.22 remain inconsistent with its printed
bonds. The amended values are consequences of actual residue graphs,
not an inference from visual similarity to another diagram.

The independent source audit also records order inconsistencies in
the displayed parabolic, Borel or kernel names of 3.1, 3.4, 3.5, 3.12,
3.13, 3.15, 3.16, 3.20, 4.24 and 4.28. We retain those discrepancies.
This theorem does not supply the abstract names that should replace
them, nor prove the remaining names merely because their orders agree.
The source's boxed kernel and generic extension dot keep their own
meaning; the dot alone is not treated as an assertion of nonsplitting.
The Aut and Cor values left blank in 3.10 are still recorded as blank.

\paragraph{Truncations.}
Restricting a coset geometry to a subset $J$ of its types retains the
same cosets and the same incidence relation. Each previously reviewed
positive truncation certificate gives an automorphism $\alpha$ of $G$
and a type permutation carrying the retained parabolics to those of
the target. The maps
\[
             G_i x\longmapsto \alpha(G_i)\alpha(x)
\]
are well defined bijections. Two cosets meet if and only if their
images meet, so these are incidence isomorphisms. Composing the
already verified certificates with the coordinate list above proves
all thirteen printed links. The record gives the exact retained source
nodes for each link. No new transporter search is needed. The eight
printed doubling assertions are separate obligations.

\paragraph{Petersen and affine residues.}
Take as Petersen vertices the ten two-element subsets of a five-element
set, with adjacency defined by disjointness. The incidence structure
has these ten vertices as points and its fifteen graph edges as lines.
This is the standard Petersen model used in the thirteen certificates.

For the affine model, set $F=\mathbb F_2[t]/(t^2+t+1)$. The polynomial
has no root in $\mathbb F_2$, so $F$ is a field of order four.
The points are $F^2$; its twenty lines are $y=ax+b$ and $x=b$.
Two distinct points lie on a unique such line, by division of their
coordinate difference when the first coordinates differ, and by the
vertical line otherwise. Thus this is $\mathrm{AG}(2,4)$.

For each of the fifteen source occurrences the certificate supplies a
bijection from the full displayed model to the actual residue graph.
For every vertex, the image of its entire neighbor set equals the
entire neighbor set of its image. This proves incidence preservation
in both directions. The model points are also checked to occupy one
entire actual type. Consequently the record gives the exact point
polarity, while retaining every printed star; no unstated convention
about reading stars on sloping diagram bonds is imposed.

\paragraph{Evidence and scope.}
The source transcription was checked on all 23 native pages by two
readers. The new coordinate calculation uses the accepted complete
finite records. Graph software only proposed the fifteen small positive
maps; their verification uses literal adjacency lists and the explicit
models above. Damaged map, roster and coordinate controls accompany
ordinary and optimized executions. The finite inventory and each
acquired review are fixed by exact file hashes. This component contributes
no new census, whole-original completion, conjecture resolution,
formalization or compression measurement. The source transcription,
earlier classification and earlier local-data proofs retain their
separate authorship and review records.
\end{document}
